% Chapitre 8 : L'Analyse des Risques

\section{Chapitre 8 : L'Analyse des Risques}



L'analyse des risques est une composante essentielle de la conduite de projet. Elle vise à identifier, évaluer et traiter les événements incertains qui pourraient affecter négativement les objectifs du projet.



\subsection{Identification des Risques}



L'identification des risques a été réalisée par une analyse systématique des dimensions techniques, sécuritaires et opérationnelles du projet. Trois catégories principales de risques ont été identifiées.



\subsubsection{Risques Techniques}



Le risque technique majeur réside dans la complexité de Next.js 16.2.4 avec le App Router. Ce framework, bien que puissant, introduit des breaking changes significatifs par rapport aux versions précédentes. La documentation évolue rapidement, et certaines fonctionnalités (Server Actions, Parallel Routes, Intercepting Routes) sont encore relativement nouvelles et peu documentées.



Un risque secondaire concerne la stack Drizzle ORM 0.45.2 + PostgreSQL. Bien que mature, la configuration des relations et des indexes nécessite une expertise SQL approfondie pour garantir les performances à grande échelle.



\subsubsection{Risques de Sécurité}



Les risques de sécurité concernent la protection des données utilisateurs et la prévention des attaques. L'injection SQL est mitigée par l'utilisation de Drizzle ORM avec des requêtes paramétrées, mais les Server Actions exposent des surfaces d'attaque potentielles si les validations côté serveur sont insuffisantes.



Les risques liés à l'authentification incluent la fuite de tokens JWT, la compromission des sessions, et les attaques par force brute. Better Auth 1.6.9 gère nativement plusieurs de ces menaces (expiration des tokens, rotation, rate limiting), mais la configuration incorrecte reste un risque.



\subsubsection{Risques de Performance}



Les risques de performance se manifestent à plusieurs niveaux. Le marketplace doit supporter un grand nombre de produits sans dégradation du temps de réponse, ce qui nécessite l'optimisation des indexes PostgreSQL et potentiellement l'ajout de pagination côté serveur. La messagerie temps réel via Redis Pub/Sub doit supporter un volume croissant de messages simultanés sans latence perceptible.



\subsection{Analyse et Stratégies de Mitigation}



Pour chaque risque identifié, une stratégie de mitigation a été définie et mise en Å“uvre.



\subsubsection{Mitigation des Risques Techniques}



La mitigation des risques techniques repose sur trois piliers. Premièrement, la veille technologique : le fichier AGENTS.md et la documentation officielle de Next.js sont consultés régulièrement pour anticiper les breaking changes. Deuxièmement, les tests end-to-end Playwright 1.59.1 détectent les régressions fonctionnelles avant qu'elles n'atteignent la branche master. Troisièmement, la modularité de l'architecture facilite le remplacement d'une technologie si elle devient obsolète.



\subsubsection{Mitigation des Risques de Sécurité}



La sécurité est assurée par plusieurs mécanismes de défense en profondeur. L'authentification Better Auth gère nativement la sécurité des sessions (tokens JWT, expiration, rotation). Les Server Actions valident systématiquement les entrées utilisateur via Zod 4.4.3. Les indexes PostgreSQL accélèrent les requêtes fréquentes (recherche par email, par slug, par catégorie) tout en limitant la surface d'attaque par injection.



\subsubsection{Mitigation des Risques de Performance}



La performance est optimisée par plusieurs techniques. Les indexes PostgreSQL sont définis dans le schéma Drizzle pour accélérer les requêtes fréquentes (recherche par email, par slug, par catégorie). Le cache Redis réduit la charge sur la base de données pour les données fréquemment accédées (sessions, profils, listes de catégories). La pagination côté serveur est implémentée pour le marketplace et la messagerie, évitant le chargement de volumes excessifs de données.



\subsection{Plan de Continuité}



En cas d'incident majeur, le plan de continuité prévoit plusieurs mesures. La restauration depuis les backups PostgreSQL est la première ligne de défense. Le système de migrations Drizzle permet de reconstruire le schéma de base de données à partir de zéro. Le versioning Git (branche master + tags) permet de revenir à une version stable du code en cas de déploiement problématique. La documentation technique (README.md, AGENTS.md) facilite l'onboarding d'un nouveau développeur en cas d'indisponibilité d'un membre de l'équipe.

